%; whizzy chapter
% -initex iniptex -latex platex -format platex -bibtex jbibtex -fmt fmt
% 以上 whizzytex を使用する場合の設定。

%     Tokyo Debian Meeting resources
%     Copyright (C) 2011 Junichi Uekawa

%     This program is free software; you can redistribute it and/or modify
%     it under the terms of the GNU General Public License as published by
%     the Free Software Foundation; either version 2 of the License, or
%     (at your option) any later version.

%     This program is distributed in the hope that it will be useful,
%     but WITHOUT ANY WARRANTY; without even the implied warranty of
%     MERCHANTABILITY or FITNESS FOR A PARTICULAR PURPOSE.  See the
%     GNU General Public License for more details.

%     You should have received a copy of the GNU General Public License
%     along with this program; if not, write to the Free Software
%     Foundation, Inc., 51 Franklin St, Fifth Floor, Boston, MA  02110-1301 USA

%  preview (shell-command (concat "evince " (replace-regexp-in-string "tex$" "pdf"(buffer-file-name)) "&"))
% 画像ファイルを処理するためにはebbを利用してboundingboxを作成。
%(shell-command "cd image201101; ebb *.png")

%%ここからヘッダ開始。

\documentclass[mingoth,a4paper]{jsarticle}
\usepackage{monthlyreport}

% 日付を定義する、毎月変わります。
\newcommand{\debmtgyear}{2011}
\newcommand{\debmtgmonth}{3}
\newcommand{\debmtgdate}{19}
% (+ (* (- 2010 2005) 12) 10) started from zero
\newcommand{\debmtgnumber}{74}

\begin{document}

\begin{titlepage}
\thispagestyle{empty}
% タイトルページ:編集必要な部分は最初のマクロに飛ばすこと

\vspace*{-2cm}
第\debmtgnumber{}回 東京エリア Debian 勉強会資料\\
\hspace*{-2cm}
\includegraphics[width=210mm]{image201003/debsen.eps}\\
\hfill{}\debmtgyear{}年\debmtgmonth{}月\debmtgdate{}日

% ここはアップデートすること
\rotatebox{10}{\fontsize{32}{32} {\gt 特集1: xxx}}

\rotatebox{10}{\fontsize{32}{32} {\gt 特集2: yyy}}

\vspace*{-2cm}
\hfill{}\includegraphics[height=6cm]{image200502/openlogo-nd.eps}
\end{titlepage}

\dancersection{Introduction}{上川 純一}

\begin{multicols}{2}
 

 今月のDebian勉強会へようこそ。これからDebianの世界にあしを踏み入れると
 いう方も、すでにどっぷりとつかっているという方も、月に一回Debianについ
 て語りませんか？

 Debian勉強会の目的は下記です。

 \begin{itemize}
 \item \underline{Debian Developer} (開発者)の育成。
 \item 日本語での「\underline{開発に関する情報}」を整理してまとめ、アップデートする。
 \item \underline{場}の提供。
 \begin{itemize}
  \item 普段ばらばらな場所にいる人々が face-to-face で出会える場を提供
	する。
  \item Debian のためになることを語る場を提供する。
  \item Debianについて語る場を提供する。
 \end{itemize}
 \end{itemize}		

 Debianの勉強会ということで究極的には参加者全員がDebian Packageをがりがり
 と作るスーパーハッカーになった姿を妄想しています。情報の共有・活用を通し
 て Debianの今後の能動的な展開への土台として、「場」としての空間を提供す
 るのが目的です。

\end{multicols}

\newpage

\begin{minipage}[b]{0.2\hsize}
 \definecolor{titleback}{gray}{0.9}
 \colorbox{titleback}{\rotatebox{90}{\fontsize{80}{80} {\gt デビアン勉強会} }}
\end{minipage}
\begin{minipage}[b]{0.8\hsize}
\hrule
\vspace{2mm}
\hrule
\begin{multicols}{2}
\tableofcontents
\end{multicols}
\vspace{2mm}
\hrule
\end{minipage}

\dancersection{事前課題}{上川 純一}

今回の事前課題は以下です:
\begin{enumerate}
 \item xxx
\end{enumerate}
この課題に対して提出いただいた内容は以下です。
\begin{multicols}{2}
{\small
% %; whizzy-master ../debianmeetingresume201101.tex
% $B0J>e$N@_Dj$r$7$F$$$k$?$a!"$3$N%U%!%$%k$G(B M-x whizzytex $B$9$k$H!"(Bwhizzytex$B$,MxMQ$G$-$^$9!#(B

\begin{prework}{ $B%-%?%O%i(B }

$B<B2H$K%j%j!<%9A0$N(BSqueeze$B$rCV$$$F$-$?$,!"FC$KLdBj$J$/;HMQ$5$l$F$$$kMM;R!#(B
$B2?$b@_Dj$7$J$/$F$b!"%5%9%Z%s%I!&%O%$%P%M!<%7%g%s$G$-$?$h$&$J!"%W%j%s%?$N@_Dj$K>/!9G:$s$@5-21$,$"$k$,!">\:Y$OK:$l$?!#(B
$B@N$KHf$Y$k$H3JCJ$K3Z$K$J$j$^$7$?!#(B
\end{prework}

\begin{prework}{ $B5HLn(B(yy\_y\_ja\_jp) }

poppler-data,cmap-adobe-*,gs-cjk-resource $B$,(B main $B$KF~$C$?$3$H$G$9$M!$(BThanks to Adobe$B!%(B
\end{prework}

\begin{prework}{ d123.itouch.com(dai) }

 \begin{itemize}
  \item $B1'Ch$,%F!<%^$G$+$C$3$$$$!*(B
  \item $B?7$7$/(BDebian$B$r3P$($k$3$H$K$J$C$?$N$G!"CzEY$$$$5!2q$G$9!#(BDebian$B=i(B
	$B?4<T!*0lG/@8$G$9!*(B
 \end{itemize}

 2$BG/$[$IA0$KEl5~$NK?@lLg3X9;$G3+:E$5$l$?5.2q$N%V!<%9$G$4CzG+$K(B
 $B$4>R2p$$$?$@$$$?$3$H$,$"$j$^$9!#$9$P$i$7$$2q$r3+:E$$$?$@$-(B
 $B$"$j$,$H$&$4$6$$$^$9!#:#2s=i$N;22C$G$9!#$I$&$>$h$m$7$/$*4j$$$7$^$9!#(B
\end{prework}

\begin{prework}{ hakase }

$B$9$$$^$;$s!"(BSqueeze$B$NL>A0;O$a$FJ9$-$^$7$?!#(B
$BJY6/$7$^$9!#(B
\end{prework}

\begin{prework}{ henrich }

$B$9$$$^$;$s!"FCCJ8D?ME*$K$&$l$7$$$H$3$m$,8+Ev$?$j$^$;$s$G$7$?!D(B
\end{prework}

\begin{prework}{ $B;3ED(B }

\preworksection{$B$&$l$7$$$3$H(B}
\begin{itemize}
 \item /bin/sh$B$,(Bdash$B$K$J$C$?!*%9%/%j%W%H5/F0$,L@$i$+$K7Z$$(B
 \item $B0MB84X78%Y!<%9$G5/F0=hM}$,$5$l$k$h$&$K$J$C$?!JAa$/$J$C$?!K(B
 \item $B%Q%C%1!<%8$,$5$i$K=<<B$7$?!J5$$E$$$F$J$+$C$?$@$1$+$b$7$l$J$$$b$N$N!"!V$3$s$J$N$^$GF~$C$F$k!WEY$,>e>:!K(B
 \item $B%j%j!<%9%9%1%8%e!<%k$,$7$C$+$j$7$F$-$?(B
\end{itemize}

\preworksection{$BJQ$o$C$?$3$H!J5$$K$J$k$3$H!K(B}
\begin{itemize}
 \item kFreeBSD$B$,F~$C$?!*$9$0;H$&$o$1$G$O$J$$$1$I!"5$$K$J$k(B
 \item arm(eb)$B$,Mn$C$3$A$?$1$I!"@$$NCfBgBN(Barmel$B$@$+$i(BOK$B$J$N$+$J!)(B
 \item Debian Live$B!#(BLiveCD$B$+$i(Bnetboot$B$^$G%+%P!<$H$$$&$3$H$GMW%A%'%C%/(B
\end{itemize}
$B%j%j!<%9$NA0$+$i<B:]$K$O;H$($F$$$?$N$G!V@hF|$+$i$$$-$J$j$&$l$7$/$J$C$?!W;v$O<B$O$J$$$b$N$N!"Ce<B$J%j%j!<%9$H$$$&46A[$G$9!#(B
\end{prework}

\begin{prework}{ dictoss($B?yK\!!E5=<(B) }

$B5/F0$,Aa$/$J$C$?!#(B
Debian GNU/kFreeBSD$B$NL>$,9-$^$C$F$-$?!#(B
\end{prework}

\begin{prework}{ yamamoto }

\preworksection{$B$&$l$7$$$3$H(B}
Squeeze $B$@$+$i!"$H$$$&$o$1$G$O$J$$$G$9$,!"(Bdebian-multimedia $B$N(B mythtv $B$,(B
 0.24 $B7O$K$J$C$F!"(Bmythtv-backend $B$,FMA3;`$9$k$3$H$,L5$/$J$j$^$7$?!#(B

\preworksection{$BJQ$o$C$?$3$H(B}
KDE $B$r>oMQ$7$F$^$9$,!"(BKlipper $B$N5sF0$,JQ$o$C$?!#$J$s$+;~!9!"$o$6$o$6(B Klipper $B$N%A%'%C%/$r$D$1$J$$$H!"Cf%\%?%s$GA*BrJ8;zNs$r%Z!<%9%H$G$-$J$$!#(B
$B860xD4::$O$^$@$@$1$I!D!#(B
$B!V%/%j%C%W%\!<%I$N%"%/%7%g%s$rM-8z$K$9$k!W$"$?$j$,$"$d$7$$(B!?
\end{prework}

\begin{prework}{ $BLnEg!!5.1Q(B }

$B$&$O$C!"IaCJ$+$i(Bunstable(sid)$B$7$+;H$C$F$J$$(B...

$B<h$j5^$.(Blenny$B$H$N0c$$$G4r$7$$;v!#(B
\begin{enumerate}
 \item kernel$B$,(B2.6.32$B7ONs$K$J$C$?!#(B(kvm$B$N%P!<%8%g%s$,>e$,$C$?!*(B)
 \item libvirt$B$,(B0.8.3$B$K$J$C$?!#!J$$$m$s$J%Q%i%a%?$,$$$m$$$m$D$1J|Bj$K(B...)
 \item gnome$B$,(B2.30$B7ONs$K$J$C$?!#(Bfreedesktop-sound-theme$B$,I8=`AuHw$5$l$?$N$G2;$,$K$.$d$+$K$J$C$?!#(B
 \item upstart$B$,Ek:\$5$l$?!#(B(sid$B;H$C$F$$$k$N$G$9$,!"(Bupstart$B$K$J$C$F$J$/(B
       $B$F$,$C$/$j!#:#$+$iJQ$($F$_$^$9(B...)
\end{enumerate}
\end{prework}

\begin{prework}{ taitioooo }

$B=i$a$F(BDebian$B$K?($C$?$N$,(Bsqueeze$B$G$7$?!#(B
$B:#$^$G(BCentos$B;H$C$F$$$^$7$?!#(B
$B=i$a$F$N0u>]$O!"!V$"!"%m%1%C%H$,$"$k!#!W$G$7$?!#(B

$B$"$H!"(Bapt-get$B$N@_Dj$,4JC1$G$h$+$C$?$G$9!#!J(Byum$B$HHf$Y$F$G$9$,!K(B
\end{prework}

\begin{prework}{ Hasegawa }

$B%$%s%9%H!<%k$G(Broot$B$G%m%0%$%s$9$k$+J9$$$F$-$?$j(B
$B@_Dj9`L\$,A}$($?$h$&$@$7!"$+$J$jJQ$o$C$?0u>]$r;}$C$?(B
64bit/32bit$B$,A*Br$G$-$k%$%s%9%H!<%i!<$K$J$C$?(B

\end{prework}

\begin{prework}{ ayakomuro }

$B$^$@%"%C%W%0%l!<%I$7$F$J$$$N$G!"=5Kv$K%"%C%W%0%l!<%I$7$?$i>\:Y=q$-$^$9(Bm(..)m
\end{prework}

\begin{prework}{$B$^$($@$3$&$X$$(B}

\preworksection{$B$&$l$7$$$3$H(B}
\begin{itemize}
 \item Squeeze$B$,%j%j!<%9$5$l$?$3$H$G(BSid$B$N3+H/$,:F3+$5$l$?$N$G!"(BCouchApp$B$d(B
       xserver-xorg-input-multitouch$B$,(BDebian$B%Q%C%1!<%8$K$J$C$?!#(B
 \item Sarge$B$G2TF/Cf$NAH9g%5!<%P$r%O!<%I%&%'%"!u%7%9%F%`O75`2=$K$h$j!"(B
       KVM/Squeeze/HP MicroServer$B$G9=C[Cf$J$N$G!"%j%j!<%9$5$l$F$A$g$&$I(B
       $BNI$+$C$?!#(B
 \item $B:#$^$G(BSid$B$G$7$+;H$($J$+$C$?%Q%C%1!<%8$,(Bstable$B$G;H$($k$h$&$K$J$C$?(B
       $B$b$N$,A}$($?!#(B(lxc, scala, xz-utils, libapache-mod-security($B$3$l$OI|3h$7(B
       $B$?!"$H$$$&$N$,@5$7$$$+(B))
 \item $B5/F0%9%/%j%W%H$,(Bbash$B$+$i(Bdash$B$KJQ$o$C$?!#(BVM$B$G$b5/F0$,B.$/$J$C$?5$(B
       $B$,$9$k!#(B($BBN46B.EY!#(B)
\end{itemize}
\preworksection{$BJQ$o$C$?$3$H(B}
\begin{itemize}
 \item libvirt$B$,(B0.8.3$B$K$J$C$?!#(Bnwfilter$B$N5!G=$,F~$C$?(B
       (/etc/libvirt/nwfilter/*.xml)$B!#(BNetfilter$B$N%k!<%k$O$9$Y$F<+J,$G=q(B
       $B$$$F$$$k$N$G!"<+J,$K$OI,MW$J$$$1$l$I!"IaCJ=q$+$J$$(BNetfilter$B$N%k!<(B
       $B%k$r=q$+$J$$?M$K<h$C$F$ONI$$$N$G$O!#(B
 \item $B%/%i%&%I4XO"$N%D!<%k$,%Q%C%1!<%8$GA}$($?!#;~Be$G$9$M!<!#(B
\end{itemize}
\end{prework}



}
\end{multicols}

\dancersection{最近のDebian関連のミーティング報告}{上川純一}
\subsection{東京エリアDebian勉強会73回目報告}

% (query-replace-regexp "<.*?>" "")
% (query-replace-regexp "^[	 ]\+" "")



\dancersection{Debian Trivia Quiz}{上川 純一}

ところで、みなさん Debian 関連の話題においついていますか？Debian関連の話
題はメーリングリストをよんでいると追跡できます。ただよんでいるだけではは
りあいがないので、理解度のテストをします。特に一人だけでは意味がわからな
いところもあるかも知れません。みんなで一緒に読んでみましょう。

今回の出題範囲は\url{debian-devel-announce@lists.debian.org} や \url{debian-devel@lists.debian.org}に投稿された
内容とDebian Project Newsからです。

\begin{multicols}{2}
% %; whizzy-master ../debianmeetingresume201101.tex
% $B0J>e$N@_Dj$r$7$F$$$k$?$a!"$3$N%U%!%$%k$G(B M-x whizzytex $B$9$k$H!"(Bwhizzytex$B$,MxMQ$G$-$^$9!#(B
%
% $B$A$J$_$K!"%/%$%:$OJL%V%i%s%A$G:n@.$7!"$N$A$K%^!<%8$7$^$9!#5U$K%^!<%8$7(B
% $B$J$$$h$&$K$7$^$7$g$&!#(B
% (shell-command "git checkout quiz-prepare")

\santaku
{RC$B%P%0$N8=>u$O$O$I$3$G3NG'$G$-$k$+(B}
{\url{http://bugs.debian.org/release-critical/}}
{\url{http://localhost/}}
{\url{http://debianmeeting.appspot.com/}}
{A}

\santaku
{Debian$BJY6/2qM=Ls%7%9%F%`$N(BURL$B$O$I$l$+(B}
{\url{http://www.2ch.net/}}
{\url{http://atnd.org/events/}}
{\url{http://debianmeeting.appspot.com/}}
{C}

\santaku
{events@debian.org$B$O$I$3$HE}9g$5$l$?$+(B}
{merchants@debian.org}
{hoge@debian.org}
{fuga@debin.org}
{A}

\santaku
{antiharassment@debian.org $B$N$&$i$K$$$J$$$N$OC/$+(B}
{Amaya Rodrigo Sastre}
{Patty Langasek}
{Kouhei Maeda}
{C}

\santaku
{Sprint$B$G$O$J$$$N$O$I$l$+(B}
{-www sprint}
{security sprint}
{tokyo sprint}
{C}

\santaku
{DACA$B$O$I$3$r8+$l$P$h$$$+(B}
{\url{http://qa.debian.org/daca/}}
{\url{http://daca.debian.org/}}
{\url{file:/tmp}}
{A}

\santaku
{DEP $B$O2?$NN,$+(B}
{Debian Enhancement Proposal}
{Device Enhancement Protocol}
{$B$G$C$W(B}
{A}

\santaku
{DEP5$B$GDs0F$5$l$F$$$k(Bdebian/copyright$B$N5!3#2DFI7A<0$O$I$&$$$&$b$N$+(B}
{S$B<0(B}
{RFC822$BIw(B}
{XML}
{B}


\end{multicols}

%-------------------------------------------------------------------------------
\dancersection{hoge}{fuga}
%-------------------------------------------------------------------------------
\index{fuu}



\printindex

\cleartooddpage

\vspace*{15cm}
\hrule
\vspace{2mm}
\includegraphics[width=2cm]{image200502/openlogo-nd.eps}
\noindent \Large \bf Debian 勉強会資料\\
\noindent \normalfont \debmtgyear{}年\debmtgmonth{}月\debmtgdate{}日 \hspace{5mm}  初版第1刷発行\\
\noindent \normalfont 東京エリア Debian 勉強会 （編集・印刷・発行）\\
\hrule

\end{document}
